\section*{Bab \ref{ch:g7:earthquakes} --- Gempa Bumi}

\begin{solution}{exo:g7:earthquakes:1}
(a) \omterm{def:g7:earthquakes:focus}{hiposentrum}; (b) \omterm{def:g7:earthquakes:intensity}{intensitas}; (c) \omterm{def:g7:earthquakes:fault}{sesar}; (d) \omterm{def:g7:earthquakes:focus}{episentrum}; (e) \omterm{def:g7:earthquakes:magnitude}{magnitudo};
(f) \omterm{def:g7:earthquakes:seismometer}{seismometer}.
\end{solution}

\begin{solution}{exo:g7:earthquakes:2}
(a) Di desa itu sendiri: \omterm{def:g7:earthquakes:focus}{episentrum} adalah titik di permukaan yang terletak
tepat di atas \omterm{def:g7:earthquakes:focus}{hiposentrum}. (b) \qty{30}{km}, lurus ke bawah. (c) Biasanya
lebih rendah: guncangan melemah makin jauh dari sumbernya. (\omterm{def:g7:earthquakes:intensity}{Intensitas} juga
bergantung pada tanahnya: di atas \omterm{def:g3:stones-pebbles-sand:sand}{lumpur} lunak, desa kedua itu bisa saja
tetap berguncang keras.)
\end{solution}

\begin{solution}{exo:g7:earthquakes:3}
(a) \omterm{def:g7:earthquakes:magnitude}{Magnitudo} (bilangan dengan satu tempat desimal). (b) \omterm{def:g7:earthquakes:intensity}{Intensitas} (apa
yang terlihat di satu tempat). (c) \omterm{def:g7:earthquakes:intensity}{Intensitas} (angka Romawi, berbeda di dua
tempat). (d) \omterm{def:g7:earthquakes:magnitude}{Magnitudo}: energi yang dilepaskan, satu langkah \omterm{def:g7:earthquakes:magnitude}{magnitudo}
lebih besar.
\end{solution}

\begin{solution}{exo:g7:earthquakes:4}
(a) Salah: para ilmuwan tidak dapat meramalkan \omterm{def:g4:volcanoes-earthquakes-first-look:earthquake}{gempa bumi}; mereka hanya
mengetahui \omterm{def:g7:earthquakes:fault}{sesar} mana yang suatu hari akan patah lagi. (b) Salah: di
\omterm{def:g2:water-on-earth:ocean}{samudra} lepas, tinggi \omterm{def:g7:earthquakes:tsunami}{tsunami} jarang lebih dari satu meter; \omterm{def:g7:earthquakes:tsunami}{tsunami} itu
makin tinggi di dekat pantai. (c) Benar. (d) Benar: 92~\% \omterm{def:g4:volcanoes-earthquakes-first-look:earthquake}{gempa bumi}
\omterm{def:g7:earthquakes:magnitude}{bermagnitudo}~6 atau lebih terjadi kurang dari \qty{300}{km} dari \omterm{def:g7:plate-tectonics:plate-boundary}{batas
lempeng}.
\end{solution}

\begin{solution}{exo:g7:earthquakes:5}
(a) Kira-kira 127 setahun yang \omterm{def:g7:earthquakes:magnitude}{bermagnitudo} 6 sampai 6.9; kira-kira 13 yang
\omterm{def:g7:earthquakes:magnitude}{bermagnitudo} 7 sampai 7.9. (b) $126.7 \div 13.12 \approx 9.7$: kira-kira
10~kali lebih sedikit. (c) $13.12 + 0.76 + 0.04 = 13.92$: kira-kira
\textbf{14} setahun. (d) $0.04 \times 50 = 2$: dua dalam 50~tahun (Sumatra
2004 dan Jepang 2011).
\end{solution}

\begin{solution}{exo:g7:earthquakes:6}
(a) Satu langkah: 32~kali. (b) Dua langkah: $32 \times 32 = 1024$,
kira-kira 1000~kali. (c) Setiap langkah membuat gelombangnya 10~kali lebih
besar: $10 \times 10 =
100$~kali lebih besar.
\end{solution}

\begin{solution}{exo:g7:earthquakes:7}
(a) $8 \times 5 = \qty{40}{km}$. (b) $8 \times 25 = \qty{200}{km}$. (c)
$120 \div 8 = \qty{15}{s}$. (d) Satu stasiun hanya memberikan jarak, bukan
arah: semua titik pada jarak itu membentuk lingkaran di sekeliling
stasiun, dan \omterm{def:g7:earthquakes:focus}{episentrumnya} bisa berada di mana saja pada lingkaran itu.
\end{solution}

\begin{solution}{exo:g7:earthquakes:8}
(a) V. (b) II. (c) VIII. (d) \omterm{def:g7:earthquakes:intensity}{Intensitas} berubah dari tempat ke tempat:
laporan (c) berasal dari tempat yang dekat dengan \omterm{def:g7:earthquakes:focus}{episentrum}, atau yang
dibangun di atas tanah lunak; laporan (a) dari tempat yang lebih jauh,
atau di atas \omterm{def:g3:stones-pebbles-sand:rock}{batuan} padat. \omterm{def:g4:volcanoes-earthquakes-first-look:earthquake}{Gempa bumi} itu tetap hanya mempunyai satu
\omterm{def:g7:earthquakes:magnitude}{magnitudo}.
\end{solution}

\begin{solution}{exo:g7:earthquakes:9}
(a) $7031 - 5561 = 1470$. (b) $1470 \div 7031 \approx 0.209$: kira-kira
\textbf{21~\%}. (c) \omterm{def:g7:plate-tectonics:plate}{Lempeng} yang menunjam turun jauh ke dalam \omterm{def:g6:inside-the-earth:layers}{mantel}.
\omterm{def:g7:plate-tectonics:plate}{Lempeng} itu masih cukup kaku untuk patah, sehingga \omterm{def:g4:volcanoes-earthquakes-first-look:earthquake}{gempa bumi} dapat
dimulai di dalamnya, ratusan kilometer di bawah permukaan.
\end{solution}

\begin{solution}{exo:g7:earthquakes:10}
(a) $2500 \div 800 = 3.125$~jam: 3~jam dan $0.125 \times 60 = 7.5$
menit, kira-kira 3~jam 8~menit. (b) 3~jam~8~menit $-$ 2~jam~17~menit $=$
51~menit lebih lambat. (c) Gelombang itu memerlukan waktu berjam-jam untuk
tiba, sedangkan letak \omterm{def:g4:volcanoes-earthquakes-first-look:earthquake}{gempa bumi} besar di bawah \omterm{def:g2:water-on-earth:ocean}{laut} dapat ditentukan
dalam beberapa menit: sebuah peringatan dapat mengirim orang ke tempat yang
tinggi jauh sebelum gelombang itu datang.
\end{solution}

\begin{solution}{exo:g7:earthquakes:11}
(a) $400 \div 800 = 0.5$~jam: 30~menit. (b) $20 \div 40 = 0.5$~jam:
30~menit. (c) 60~menit. (d) \qty{20}{km} terakhir memerlukan waktu yang
sama lamanya dengan \qty{400}{km} pertama: gelombang itu melambat luar
biasa di perairan dangkal. (e) Bagian depan gelombang melambat sementara
air di belakangnya masih datang dengan cepat: air itu menumpuk, dan
gelombangnya makin tinggi.
\end{solution}

\begin{solution}{exo:g7:earthquakes:12}
(a) Tiga langkah: $32 \times 32 \times 32 = 32\,768$, kira-kira
\textbf{33\,000}~kali. (b) $32\,768 \div 127 \approx 258$: kira-kira
260~tahun dari semua \omterm{def:g4:volcanoes-earthquakes-first-look:earthquake}{gempa bumi} \omterm{def:g7:earthquakes:magnitude}{bermagnitudo}~6 di seluruh dunia. (c) Hampir
seluruh energi dilepaskan oleh segelintir \omterm{def:g4:volcanoes-earthquakes-first-look:earthquake}{gempa bumi} yang paling besar:
satu gempa \omterm{def:g7:earthquakes:magnitude}{bermagnitudo}~9 saja melebihi gempa-gempa \omterm{def:g7:earthquakes:magnitude}{bermagnitudo}~6 selama
berabad-abad.
\end{solution}

\begin{solution}{pb:g7:earthquakes:1}
\textbf{1.} A: $45 - 25 = 20$~s. B: $29 - 16 = 13$~s. C: $34 - 19 =
15$~s.

\textbf{2.} Stasiun B: selisihnya paling pendek, dan makin pendek
selisihnya, makin dekat \omterm{def:g4:volcanoes-earthquakes-first-look:earthquake}{gempa buminya}. (Gelombangnya yang lebih cepat juga
tiba paling awal.)

\textbf{3.} A: $8 \times 20 = \qty{160}{km}$. B: $8 \times 13 =
\qty{104}{km}$. C: $8 \times 15 = \qty{120}{km}$.

\textbf{4.} Semua titik yang berjarak \qty{160}{km} dari A: sebuah lingkaran
berjari-jari \qty{160}{km} yang berpusat di A. \omterm{def:g7:earthquakes:focus}{Episentrum} terletak di suatu
tempat pada lingkaran itu.

\textbf{5.} A: $160 \div 20 = 8$~kotak. B: $104 \div 20 = 5.2$~kotak.
C: $120 \div 20 = 6$~kotak.

\textbf{6.} E terletak $3 + 6 = 9$~kotak ke timur dan 14.5~kotak ke utara
dari O, di \omterm{def:g2:water-on-earth:ocean}{laut} di antara ketiga stasiun; lingkaran A dan lingkaran B
melaluinya.

\textbf{7.} Di bawah \omterm{def:g2:water-on-earth:ocean}{laut}. Gerakan dasar \omterm{def:g2:water-on-earth:ocean}{laut} yang tiba-tiba dapat memicu
\omterm{def:g7:earthquakes:tsunami}{tsunami}: pantai-pantai dan pulau itu harus diberi peringatan.

\textbf{8.} Satu langkah \omterm{def:g7:earthquakes:magnitude}{magnitudo}: energinya kira-kira 32~kali lebih banyak.

\textbf{9.} $12 \times 20 = \qty{240}{km}$.

\textbf{10.} $240 \div 800 = 0.3$~jam, dan $0.3 \times 60 = 18$~menit.

\textbf{11.} 5~menit sesudah \omterm{def:g4:volcanoes-earthquakes-first-look:earthquake}{gempa bumi}.

\textbf{12.} $18 - 5 = 13$: penduduk Port Hale mempunyai waktu kira-kira
\textbf{13~menit} untuk mencapai tempat yang tinggi. (Gelombang itu melambat
di perairan dangkal di dekat pulau, sehingga ada sedikit waktu tambahan,
tetapi tidak ada yang boleh mengandalkannya: mereka harus segera pergi.)
\end{solution}
